\documentclass[11pt,a4paper]{article}

\usepackage[margin=2cm]{geometry}
\usepackage{booktabs}
\usepackage{enumitem}
\usepackage{titlesec}
\usepackage{parskip}
\usepackage{hyperref}
\usepackage{graphicx}
\usepackage{array}
\usepackage{xcolor}
\usepackage{fancyhdr}
\usepackage{float}

% ── Page style ────────────────────────────────────────────────────────────────
\pagestyle{fancy}
\fancyhf{}
\lhead{\small CSE352 Computer Graphics}
\chead{\small \textit{Final Project Report}}
\rhead{\small Spring 2026}
\cfoot{\small Page \thepage}
\renewcommand{\headrulewidth}{0.4pt}

% ── Section formatting ────────────────────────────────────────────────────────
\titleformat{\section}{\normalfont\large\bfseries}{}{0em}{}[\titlerule]
\titlespacing{\section}{0pt}{10pt}{4pt}

% ── Table column types ────────────────────────────────────────────────────────
\newcolumntype{L}[1]{>{\raggedright\arraybackslash}p{#1}}

% ─────────────────────────────────────────────────────────────────────────────
\begin{document}

% ── Title block ───────────────────────────────────────────────────────────────
\begin{center}
    {\LARGE\bfseries Charizard's Chase: A WebGL Mini-Game}\\[4pt]
    {\large Musab Suhail \quad|\quad ERP: \texttt{26923}}\\[2pt]
    {\normalsize CSE352 Computer Graphics\quad--\quad Spring 2026}\\[2pt]
    {\normalsize \textbf{Project Direction:} Mini Game (Pokémon-inspired)}
\end{center}

\vspace{-4pt}\noindent\rule{\linewidth}{0.6pt}\vspace{6pt}

% ── 1. Final Summary ──────────────────────────────────────────────────────────
\section{Final Summary}

\textit{Charizard's Chase} is a fully playable, end-to-end WebGL mini-game built
on Three.js r160 without any build tools, running directly in the browser via
ES-module importmaps. The game has a complete narrative arc: an opening cutscene
shows Charizard preparing to eat a Magikarp on a wooden table when Golbat swoops
in and steals it; the player then chases Golbat through an infinite procedural
world; catching Golbat triggers a closing cutscene in which Golbat and Magikarp
land on a second (wider) table, Charizard stands before it, and cooks the
Magikarp with a swelling flame effect.

Core gameplay features include a velocity-based Golbat AI with panic behaviour, a
coin-collection boost system, a three-heart health system with invincibility frames,
sludge-bomb projectile attacks, and a full pause/retry menu. The world uses pooled
GLB rock, cloud, and tree assets over an infinite tiling grass texture with
world-space UV anchoring, and procedural cone mountains with tapered collision
geometry. All three advanced techniques --- toon shading, post-processing bloom, and
PCF shadow mapping --- are active and visible throughout.

Since the progress report, the project received a full round of correctness fixes,
rendering-quality improvements, and CPU/GPU performance optimisations detailed in
Section~6.

% ── 2. Baseline Techniques ───────────────────────────────────────────────────
\section{Baseline Techniques Status}

The four baseline techniques listed in the project proposal are all completed.

\begin{tabular}{@{} L{3.5cm} L{2.2cm} L{8.8cm} @{}}
\toprule
\textbf{Technique} & \textbf{Status} & \textbf{Notes} \\
\midrule
Phong Lighting Model
  & Completed
  & \texttt{DirectionalLight} (sun, intensity 1.2) + \texttt{AmbientLight}
    (fill) + \texttt{PointLight} (Charizard tail flame, repositioned every
    frame). During the outro cooking scene the point light intensity ramps from
    3 to 11 and its world position tracks the Magikarp on the table, producing a
    visible heat-glow bloom pulse. \\[4pt]
Texture Mapping
  & Completed
  & Grass diffuse texture (\texttt{grass\_diff\_1k.jpg}, 1K, Polyhaven CC0)
    applied to the infinite ground plane with a repeat factor of 80 and a
    world-space UV counter-offset each frame so the texture appears anchored to
    world space rather than scrolling with the camera. GLB model diffuse
    textures are preserved through the toon-material conversion via
    \texttt{map} pass-through, satisfying the ``character surfaces'' requirement
    from the proposal. \\[4pt]
Multiple Light Sources
  & Completed
  & Three distinct types in simultaneous use: \texttt{DirectionalLight}
    (sun/shadows), \texttt{AmbientLight} (fill), \texttt{PointLight}
    (Charizard tail flame; directional fire, as proposed). Sludge bombs
    additionally carry a per-bomb purple \texttt{PointLight}
    (intensity 6, range 10). \\[4pt]
Camera System
  & Completed
  & Third-person dynamic chase camera during gameplay (lerp-smoothed, 6 units
    behind, 3 units above, with look-ahead), satisfying the proposal spec.
    Supplemented with scripted cinematic cameras for the intro (low-angle drift,
    swoop tracking) and outro (orbital descent, snap to table framing,
    pullback). \\
\bottomrule
\end{tabular}

\medskip
\noindent\textit{Note on shaders:} Although not listed as a separate baseline
item in the proposal, the Three.js rendering pipeline uses GLSL shaders
internally for all passes. All mesh materials are \texttt{MeshToonMaterial}
with \texttt{roughness=1}, \texttt{metalness=0}, and
\texttt{envMapIntensity=0}, and the post-processing stack runs
\texttt{UnrealBloomPass} as a full-screen GLSL quad. Scene-graph transformations
(YXZ Euler yaw/pitch/roll for flight, quaternion-derived movement, Golbat
procedural bank/pitch) and correct shadow normals via
\texttt{PCFSoftShadowMap} are applied throughout.

% ── 3. Advanced Techniques ───────────────────────────────────────────────────
\section{Advanced Techniques Status}

\textbf{Toon Shading (3/5) --- Completed.}
\texttt{MeshToonMaterial} is applied to all geometry in the scene. All loaded
GLB models (Charizard, Golbat, Magikarp, rocks, clouds, trees) have their
PBR materials overwritten on load via a post-load \texttt{traverse} that sets
\texttt{roughness = 1}, \texttt{metalness = 0}, and
\texttt{envMapIntensity = 0} while preserving any diffuse \texttt{map} texture,
ensuring a consistent flat, matte, cel-shaded look across both procedural and
imported assets.
\textit{Challenge:} PBR materials from GLB exports clashed with the toon
aesthetic. Resolved with the post-load traverse approach.

\medskip
\textbf{Bloom (4/5) --- Completed.}
\texttt{EffectComposer} with \texttt{UnrealBloomPass} (strength 0.8, radius 0.4,
threshold 0.85). The tail flame mesh (\texttt{MeshToonMaterial}, emissive
\texttt{0xff2200}) glows through bloom. Gold coins use emissive
\texttt{0xFFAA00}. During the outro cooking scene the \texttt{PointLight}
intensity is animated from 3 to 11 with its position locked above the Magikarp
on the table, creating a visible bloom flare as the flame cooks the Magikarp in
place.
\textit{Challenge:} Bloom overexposed toon shading at high strength. Balanced
by lowering strength and raising threshold.

\medskip
\textbf{Shadow Mapping (4/5) --- Completed.}
\texttt{PCFSoftShadowMap} with a 1024\texttimes1024 depth map from the directional
sun light; frustum manually sized (\(\pm40\) units, far 200). All characters,
Magikarp, obstacles, and the intro/outro tables cast and receive shadows; the
infinite ground plane receives them. Shadow bias has been tuned to
\texttt{sunLight.shadow.bias = -0.0001} to eliminate acne on steep mountain
slopes.

\medskip
\noindent Advanced technique effort scores: Toon Shading 3/5 + Bloom 4/5 +
Shadow Mapping 4/5 = \textbf{11/10} \checkmark
\quad(satisfies combined $\geq$ 6/10 and at least one technique $\geq$ 3/5).

% ── 4. Cutscene Implementation ───────────────────────────────────────────────
\section{Cutscene Implementation}

Both cutscenes are driven by a phase-based timeline engine: each phase object
carries \texttt{onEnter}, \texttt{onEnterCapture}, \texttt{update(t, elapsed)},
and \texttt{onExit} callbacks. \texttt{t} is the normalised phase-local progress
$[0,1]$ and \texttt{elapsed} is the running total since cutscene start. Smooth
motion uses the cubic ease \(t^2(3-2t)\).

\medskip
\textbf{Intro (6 s, 6 phases).}
The opening scene shows Charizard standing behind a wooden table on which
a Magikarp is placed. The table is constructed entirely from Three.js primitives:
a \texttt{BoxGeometry} (1.0 \texttimes{} 0.07 \texttimes{} 0.6) tabletop and four
\texttt{CylinderGeometry} legs, all using \texttt{MeshToonMaterial} in a warm
brown (\texttt{0x7a4e2d}), casting and receiving shadows. The Magikarp rests on
the table surface at the correct world height and flops using a frame-synchronised
sine animation driven by \texttt{elapsed} (replacing a previous
\texttt{Date.now()} approach that caused jitter under FPS drops). Golbat swoops
down from off-screen directly toward the table to snatch the Magikarp, which is
then carried away as Golbat escapes.

\medskip
\textbf{Outro (7 s, 6 phases).}
The victory cutscene uses a wider table (2.0 \texttimes{} 0.07 \texttimes{} 0.6),
double the intro width, to accommodate both Golbat and Magikarp simultaneously.
Golbat descends from the sky and lands on the table surface
(computed as \texttt{TABLE\_SURFACE\_Y + golbatHalfHeight} from a run-time
bounding-box measurement), with Charizard and Pikachu standing on the ground
behind it facing the camera --- mirroring the intro composition. Once Golbat
has settled, a cooking phase begins: the flame swells, a \texttt{PointLight}
blooms above the Magikarp's fixed table position, and the camera pulls back to
frame the entire scene.

% ── 5. Screenshots ───────────────────────────────────────────────────────────
\section{Screenshots / Renders}

\begin{figure}[H]
    \centering
    \includegraphics[width=0.48\linewidth]{screenshot1}
    \hfill
    \includegraphics[width=0.48\linewidth]{screenshot2}
    \\[6pt]
    \includegraphics[width=0.48\linewidth]{screenshot3}
    \hfill
    \includegraphics[width=0.48\linewidth]{screenshot4}
\end{figure}

% ── 6. Technical Improvements & Optimisations ────────────────────────────────
\section{Technical Improvements \& Optimisations}

The following improvements were made after the progress report submission.

\medskip
\textbf{Bug Fixes.}
\begin{itemize}[leftmargin=*, itemsep=2pt]
    \item \texttt{standingCharizard} was referenced in the Escape-key handler
          without being imported, causing a \texttt{ReferenceError} crash when
          pressing Escape to skip the intro. Fixed by adding it to the import.
    \item Coin respawn used \texttt{setTimeout(8000)}, whose callbacks survived
          \texttt{resetGame()} calls and fired 8~seconds after collection even
          after a retry. Replaced with a per-coin \texttt{hiddenTimer} counter
          incremented each frame; a new \texttt{resetCoins()} export resets all
          timers and visibility on retry.
    \item Magikarp flopping animations in the intro used \texttt{Date.now()}
          wall-clock time, causing visual jitter under FPS drops and
          desynchronisation from the \(t\)-based camera motion. All five
          occurrences replaced with \texttt{elapsed} (the frame-synced running
          time), which is now passed as a second argument to each phase's
          \texttt{update(t, elapsed)} callback.
\end{itemize}

\medskip
\textbf{Rendering Quality.}
\begin{itemize}[leftmargin=*, itemsep=2pt]
    \item \texttt{sunLight.shadow.bias = -0.0001} added to eliminate shadow
          acne on steep mountain faces (previously listed as remaining work).
    \item Camera far plane reduced from 400 to 250 units. Scene fog ends at 220
          units, so the previous 180-unit gap was rendering invisible geometry
          every frame.
    \item \texttt{renderer.setPixelRatio(Math.min(window.devicePixelRatio, 2))}
          added at initialisation and in the resize handler, preventing 3\(\times\)
          resolution rendering on high-DPI displays that previously halved
          framerate.
\end{itemize}

\medskip
\textbf{CPU / GC Performance.}
\begin{itemize}[leftmargin=*, itemsep=2pt]
    \item \textbf{Scratch vectors:} Six module-level \texttt{THREE.Vector3}
          instances (\texttt{\_tmpV1}--\texttt{\_tmpV6}) declared once and
          reused throughout \texttt{updateGameplay()}. Previously, nine or more
          \texttt{new THREE.Vector3()} calls per frame (fire glow, Golbat AI,
          compass, camera) produced roughly 500 short-lived heap objects per
          second, triggering frequent minor GC pauses. The \texttt{getCharizardForward()}
          helper was similarly patched to reuse a module-level scratch.
    \item \textbf{Collision loop culling:} Both the O(130) Golbat-obstacle
          avoidance loop and the O(130) Charizard-collision loop now perform a
          squared-distance fast-reject (\(\Delta x^2 + \Delta z^2 > 225\)) before
          any \texttt{sqrt} or destructuring. In typical gameplay fewer than 15
          obstacles pass the test, reducing per-frame collision work by
          approximately 88\%.
    \item \textbf{Shared sludge-bomb geometry:} \texttt{makeBomb()} previously
          allocated a new \texttt{SphereGeometry}, \texttt{TorusGeometry}, and two
          \texttt{MeshToonMaterial} instances every 4 seconds. These four objects
          are now module-level singletons shared by all active bombs.
\end{itemize}

% ── 7. Challenges Encountered ────────────────────────────────────────────────
\section{Challenges Encountered}

A forward-vector orientation bug caused Charizard to fly sideways until a
canonical \(+Z\) forward convention was established with quaternion-derived
movement. Golbat's AI initially maintained a fixed offset regardless of player
speed, making it uncatchable; replaced with a velocity-capped system so a boosted
Charizard can close the gap. Tapered bounding-box collision against tall mountains
caused phantom crashes far from the geometry; replaced with a per-type
cylinder/tapered-cone check that matches each obstacle's visual shape.

Placing cutscene models precisely on the table required computing
\texttt{golbatHalfHeight} at runtime from a \texttt{Box3} measurement
immediately after loading, since model origin points vary across GLB exports.
This measurement drives the outro landing target
(\texttt{TABLE\_SURFACE\_Y + golbatHalfHeight}) so Golbat's feet rest on the table
surface regardless of the source model's coordinate system.

Attaching the tail flame to the animated tail bone remains a positional offset
rather than a bone-parented mesh, because the GLB skeleton lives inside
\texttt{SkinnedMesh.skeleton.bones} rather than the scene graph traversal.

% ── 8. Deferred Items ────────────────────────────────────────────────────────
\section{Deferred Items}

The following items were not completed within the project timeline:

\begin{enumerate}[leftmargin=*, label=\arabic*.]
    \item \textbf{Skybox / HDRI environment} --- the proposal's Week~1 plan
          included a skybox. \texttt{kiara\_1\_dawn\_1k.hdr} (Polyhaven CC0) was
          acquired but not integrated; the scene background remains a solid sky
          colour (\texttt{0x87CEEB}) with matching linear fog to blend the
          horizon. A proper HDRI skybox would replace the solid colour and
          provide image-based ambient lighting.
    \item \textbf{Custom GLSL pass} --- a toon-outline or ramp shader was planned
          but deferred in favour of completing the cutscene narratives and
          optimisation work. The Three.js \texttt{MeshToonMaterial} and
          \texttt{UnrealBloomPass} built-in shaders fulfil the shading
          requirements.
    \item \textbf{Tail-flame bone parenting} --- the flame mesh uses a positional
          offset from Charizard's root; skeletal attachment is non-trivial
          because the GLB skeleton lives inside
          \texttt{SkinnedMesh.skeleton.bones} rather than the traversable scene
          graph.
\end{enumerate}

% ── 9. Scope Adjustments ─────────────────────────────────────────────────────
\section{Scope Adjustments}

\begin{tabular}{@{} L{4cm} L{10cm} @{}}
\toprule
\textbf{Change} & \textbf{Justification} \\
\midrule
Murkrow \(\rightarrow\) Golbat
  & A suitably rigged Murkrow GLB was unavailable. Golbat fills the same
    role and a quality rigged asset was sourced from Sketchfab. \\[4pt]
Lighter/cigarette \(\rightarrow\) Magikarp
  & Replaced the prop-theft premise with a food-theft premise for a more
    visual and humorous payoff. Golbat steals Charizard's Magikarp from a
    table; the outro shows Charizard cooking it on a second table. \\[4pt]
Pikachu easter-egg
  & Originally deferred; now fully integrated --- Pikachu appears beside
    Charizard in both cutscenes and rides on Charizard's back during
    gameplay via a dedicated \texttt{pikachu-flying-on-charizard.glb}. \\[4pt]
Codebase architecture
  & Single-file prototype split into twelve ES modules
    (\texttt{scene}, \texttt{world}, \texttt{game}, \texttt{controls},
    \texttt{charizard}, \texttt{golbat}, \texttt{magikarp}, \texttt{pikachu},
    \texttt{sludge}, \texttt{coins}, \texttt{intro}, \texttt{outro}).
    No techniques added or removed. \\[4pt]
Pikachu source: Thangs $\rightarrow$ Sketchfab
  & The Thangs asset was not suitable for the required riding animation.
    A higher-quality rigged GLB was sourced from Sketchfab instead;
    the asset is properly credited. Two Pikachu variants are used:
    a standing model for cutscenes and a
    \texttt{pikachu-flying-on-charizard.glb} riding model during gameplay. \\[4pt]
Custom GLSL pass deferred
  & Cutscene polish and performance work were prioritised. Three.js built-in
    shaders satisfy the shading requirements. \\[4pt]
Skybox deferred
  & HDRI file acquired but not wired in; solid sky colour with matching fog
    provides a visually consistent background. \\
\bottomrule
\end{tabular}

% ── 10. External Assets ───────────────────────────────────────────────────────
\section{External Assets}

\begin{tabular}{@{} L{4.2cm} L{2.8cm} L{6.5cm} @{}}
\toprule
\textbf{Asset} & \textbf{Source} & \textbf{Credit / Link} \\
\midrule
Charizard flying animation GLB
  & Sketchfab
  & \url{https://skfb.ly/prsZz} \\[4pt]
\texttt{charizard.glb} (standing)
  & Sketchfab
  & \url{https://skfb.ly/ouA78} \\[4pt]
\texttt{golbat.glb}
  & Sketchfab
  & \url{https://skfb.ly/oPXDW} \\[4pt]
\texttt{magikarp.glb}
  & Sketchfab
  & \url{https://skfb.ly/6ZXn9} \\[4pt]
\texttt{pikachu.glb}
  & Sketchfab
  & \url{https://skfb.ly/pGUWE}
    (proposed source was Thangs; replaced with Sketchfab asset) \\[4pt]
\texttt{pikachu-flying-on-charizard.glb}
  & Sketchfab
  & \url{https://skfb.ly/p8NOI} \\[4pt]
\texttt{grass\_diff\_1k.jpg}
  & Polyhaven
  & Polyhaven (CC0) \\[4pt]
\texttt{kiara\_1\_dawn\_1k.hdr}
  & Polyhaven
  & Polyhaven (CC0) --- acquired, not integrated \\[4pt]
\texttt{rock.glb, cloud.glb, tree.glb}
  & Quaternius / itch.io
  & Quaternius (CC0) \\
\bottomrule
\end{tabular}

\medskip
All Three.js library code is loaded from the \texttt{unpkg.com} CDN
(Three.js r0.160.0, MIT licence). No other third-party runtime dependencies
are used.

\end{document}
